% =============================================================================
% cap08.tex — DETALHAMENTO CONSTRUTIVO
% =============================================================================

\chapter{DETALHAMENTO CONSTRUTIVO}\label{cap:detalhamento}

\section{CRITÉRIO DE MARGEM DAS LIGAÇÕES}\label{sec:criterio-ligacoes}

As ligações recebem tratamento distinto do dispensado aos elementos. A razão
está na natureza da ruptura. Uma viga que atinge a resistência à flexão escoa
progressivamente e redistribui esforço; um grupo de pinos que provoca o
fendilhamento da madeira rompe de forma súbita, sem aviso e sem redistribuição.

Adotou-se, por isso, limite de projeto de $\eta \leq 0{,}75$ para os modos
frágeis, mantendo-se $\eta \leq 1{,}00$ para os dúcteis. A margem adicional de
vinte e cinco por cento não decorre de prescrição normativa, e sim de decisão de
projeto justificada pelo uso público da edificação. A
\autoref{tab:modos-margem} apresenta a classificação aplicada e o resultado de
cada modo, segundo a \autoref{tab:modos-ruptura}. As exceções admitidas ao
limite de 0,75 são a compressão no concreto e na madeira e o escoamento do aço,
por serem modos dúcteis com capacidade de redistribuição.

\tabela
{Classificação dos modos de ruptura das ligações e taxas obtidas}
{\footnotesize
\begin{tabular}{llccc}
\toprule
    \textbf{Modo de ruptura} & \textbf{Natureza} & \textbf{Limite} & \textbf{$\eta$} & \textbf{Situação} \\
\midrule
    Fendilhamento em L1          & Frágil & 0,75 & 0,7222 & Atende \\
    Fendilhamento em L1c         & Frágil & 0,75 & 0,7222 & Atende \\
    Pinos de L1                  & Frágil & 0,75 & 0,6634 & Atende \\
    Pinos de L1c                 & Frágil & 0,75 & 0,6634 & Atende \\
    Ruptura de bloco em L1       & Frágil & 0,75 & 0,0051 & Atende \\
    Ruptura de bloco em L1c      & Frágil & 0,75 & 0,0051 & Atende \\
    Chumbadores de L3            & Frágil & 0,75 & 0,0000 & Atende \\
    Chapa metálica de L1         & Dúctil & 1,00 & 0,0956 & Atende \\
    Chapa metálica de L1c        & Dúctil & 1,00 & 0,0956 & Atende \\
    Apoio da base central L3c    & Dúctil & 1,00 & 0,1328 & Atende \\
\bottomrule
\end{tabular}}
{Elaborado pelo autor (2026)}
{modos-margem}
{}
{}

Todos os modos frágeis situam-se abaixo de 0,75, e os dois mais solicitados são
o fendilhamento das ligações viga-pilar, com 0,7222.

\section{LIGAÇÃO ENTRE VIGA E PILAR}\label{sec:ligacao-viga-pilar}

\subsection{Concepção}\label{ssec:concepcao-l1}

A ligação entre a viga e o pilar é materializada por chapa metálica embutida em
rasgo central da peça de madeira, atravessada por pinos passantes. A solução
apresenta duas vantagens em relação a chapas externas aparafusadas: a chapa fica
protegida da intempérie pela própria madeira, o que importa em classe de serviço
3, e a ligação trabalha em corte duplo, com dois planos de corte por pino.

A \autoref{tab:ligacao-l1} apresenta o dimensionamento, conforme os itens 8.2,
8.5.1.1 e 8.6 da EN 1995-1-1 \cite{EN1995-1-1:2004}, com $\gamma_M = 1{,}30$ e
$k_{mod} = 0{,}65$. O diâmetro situa-se na faixa de 6 a 30 mm exigida pela norma,
e a chapa é em aço S275 de 25 mm de espessura.

\tabela{Dimensionamento da ligação entre viga e pilar}{\footnotesize
\begin{tabular}{lrl}
\toprule
    \textbf{Grandeza} & \textbf{Valor} & \textbf{Unidade} \\
\midrule
    Força resultante de cálculo        &  99,61 & kN \\
    Ângulo da resultante com as fibras &  84,96 & graus \\
    Diâmetro do pino                   &  24    & mm \\
    Número de pinos                    &   6    & \\
    Número de planos de corte por pino &   2    & \\
    Resistência por pino e plano       &  25,03 & kN \\
    Resistência total da ligação       & 150,15 & kN \\
    Espaçamento entre pinos            &  76,2  & mm \\
    Distância à extremidade carregada  & 168,0  & mm \\
    Comprimento ocupado pela ligação   & 240,0  & mm \\
    Altura ocupada pela ligação        & 551,6  & mm \\
    Altura disponível na seção         & 1.000  & mm \\
\bottomrule
\end{tabular}}{Elaborado pelo autor (2026)}{ligacao-l1}{}{}

\subsection{Modo de ruptura governante}\label{ssec:modo-governante-l1}

O aspecto determinante desta ligação é o ângulo da força resultante em relação
às fibras, de 84,96$^{\circ}$. Como o modelo biapoiado anula o esforço
horizontal, a resultante é praticamente vertical, isto é, quase perpendicular às
fibras da viga. Disso decorrem duas consequências.

A primeira é a redução da resistência ao embutimento. O coeficiente $k_{90}$
assume o valor 1,71, e a resistência ao embutimento cai de 23,99 MPa na direção
paralela para 14,08 MPa na direção efetiva da força. A verificação segundo a
teoria de Johansen resulta governada pelo modo que combina o embutimento da
madeira com a formação de duas rótulas plásticas no pino, conduzindo a 25,03 kN
por pino e plano de corte, contra 50,00 kN do modo de embutimento puro.

A segunda, e mais relevante, é o fendilhamento. Uma força perpendicular às
fibras tende a separá-las, e essa é a solicitação que o item 8.1.4 da
EN 1995-1-1 \cite{EN1995-1-1:2004} verifica. Com a força de 99 kN aplicada a
altura efetiva de 680 mm, correspondente a 0,68 da altura da seção, a
resistência ao fendilhamento da madeira não reforçada é de 103 kN, do que
resulta taxa de 0,961.

Esse valor atende ao limite normativo, porém excede em muito o limite de projeto
de 0,75 estabelecido para modos frágeis. Adotou-se, por isso, reforço
perpendicular às fibras conforme o item 8.7.2 da mesma norma, com oito
parafusos autoperfurantes de 10 mm e comprimento efetivo de 400 mm, dispostos na
região da ligação. Com o reforço, a taxa cai para 0,7222, governada pelo
arrancamento dos parafusos, com 0,72, e não pela resistência do seu aço, com
0,50.

A ligação da viga sobre o pilar central, designada L1c, difere apenas por
empregar duas chapas, uma para cada viga, e apresenta resultados idênticos.

\subsection{Modos verificados e não governantes}\label{ssec:modos-nao-governantes}

A ruptura de bloco, verificada conforme o anexo A, apresenta taxa de 0,0051. A
resistência característica é de 3.318 kN, governada pela parcela de tração, mais
de trinta vezes superior à solicitação. O arranjo em linha única de seis pinos,
com apenas um pino por fileira paralela às fibras, é o responsável por essa
folga: não há volume de madeira delimitado por um grupo capaz de se destacar.

Esse mesmo arranjo explica o número efetivo de pinos. Com uma única coluna de
pinos na direção das fibras, o coeficiente de redução por fileira vale a
unidade, e o número efetivo iguala o número real, seis. Fosse o arranjo em
matriz, a redução seria significativa.

A chapa metálica, verificada conforme a EN 1993-1-8 \cite{EN1993-1-8:2005},
apresenta taxa de 0,0956, governada pelo cisalhamento do pino. As verificações
de pressão de contato, de seção líquida e de solda resultam em 0,040, 0,021 e
0,048. A espessura de 25 mm não decorre dessas verificações, e sim da exigência
de que a chapa central seja rígida o bastante para que o pino forme rótula
plástica, o que requer espessura não inferior ao seu diâmetro.

O deslizamento da ligação em serviço e a alternância de esforços foram
verificados conforme os itens 2.4, 7.1 e 8.1.5, resultando em taxas de 0,0472 e
0,7333.

\section{LIGAÇÃO ENTRE PILAR E FUNDAÇÃO}\label{sec:ligacao-base}

A base dos pilares emprega chapa metálica soldada a chumbadores. Nas bases de
extremidade, designadas L3, o esforço horizontal transmitido é nulo, em
consequência do modelo adotado, e os quatro chumbadores correspondem a mínimo
construtivo, com taxa nula. Nas bases centrais, designadas L3c, o esforço normal
de 199,2 kN produz tensão de contato de 2,213 MPa no concreto, com taxa de
0,1328.

As chapas de base medem 400 por 400 mm. A espessura de 8 mm nas bases de
extremidade corresponde ao mínimo construtivo; nas centrais, a verificação da
flexão do balanço da chapa exige 10,1 mm, adotando-se 12,5 mm por
disponibilidade comercial.

\section{ANCORAGEM AO BLOCO DE FUNDAÇÃO}\label{sec:ancoragem}

\subsection{Esforços de arrancamento}\label{ssec:arrancamento}

A sucção do vento inverte a reação vertical nas bases, e o arrancamento
resultante depende da posição da base na edificação. A
\autoref{tab:arrancamento} apresenta os valores, na combinação de sucção da
EN 1990 \cite{EN1990:2002}. A base de canto acumula o arrancamento de sucção, de
83,9 kN, com o esforço descendente dos dois painéis de contraventamento, de
193,6 kN, invertido sob a mesma combinação.

\tabela{Esforço de arrancamento por posição de base}{\footnotesize
\begin{tabular}{lcrc}
\toprule
    \textbf{Posição da base} & \textbf{Quantidade} & \textbf{Arrancamento (kN)} & \textbf{$\eta$} \\
\midrule
    Extremidade & 14 &  83,9 & 0,262 \\
    Central     &  9 & 151,8 & 0,474 \\
    Canto       &  4 & 277,5 & 0,578 \\
\bottomrule
\end{tabular}}{Elaborado pelo autor (2026)}{arrancamento}{}{}

A base de canto é a mais solicitada, por acumular efeitos que não coincidem nas
demais, e é ela que determina a ancoragem.

\subsection{Verificação da ancoragem de canto}\label{ssec:ancoragem-canto}

Adotaram-se seis chumbadores moldados com cabeça, de 24 mm e classe 8.8. A
verificação divide-se em duas partes independentes: a resistência do aço do
chumbador, segundo a EN 1993-1-8 \cite{EN1993-1-8:2005}, e a resistência do
concreto que o envolve, segundo a EN 1992-4 \cite{EN1992-4:2018}. A
\autoref{tab:ancoragem-canto} apresenta ambas. O aço segue a tabela 3.4 e o item
3.4 da EN 1993-1-8 \cite{EN1993-1-8:2005}, com $\gamma_{Ms} = 1{,}5$; o concreto
C25/30 segue as fórmulas (7.1) a (7.12) da EN 1992-4 \cite{EN1992-4:2018}, com
profundidade efetiva de 400 mm, afastamento de borda de 600 mm e espaçamento de
400 mm; e o esmagamento sob a chaveta segue o item 6.7 da
EN 1992-1-1 \cite{EN1992-1-1:2004}, para áreas parcialmente carregadas.

\tabela{Verificação da ancoragem da base de canto}{\footnotesize
\begin{tabular}{lrlc}
\toprule
    \textbf{Grandeza} & \textbf{Valor} & \textbf{Unidade} & \textbf{$\eta$} \\
\midrule
    \multicolumn{4}{l}{\textit{Resistência do aço}} \\
    Tração solicitante          &   277,48 & kN & \\
    Tração resistente do grupo  & 1.016,64 & kN & 0,273 \\
    Cortante solicitante        &   159,01 & kN & \\
    Cortante nos chumbadores    &     0,00 & kN & 0,000 \\
    Interação                   &          &    & 0,195 \\
\addlinespace
    \multicolumn{4}{l}{\textit{Resistência do concreto}} \\
    Cone de arrancamento        &   527,41 & kN & 0,526 \\
    Extração pela cabeça        & 1.781    & kN & \\
    Fendilhamento               &   527,41 & kN & 0,526 \\
    Esmagamento sob a chaveta   &   266,67 & kN & 0,596 \\
\bottomrule
\end{tabular}}{Elaborado pelo autor (2026)}{ancoragem-canto}{}{}

O cortante de 159,01 kN não solicita os chumbadores. Ele é transmitido ao bloco
por chaveta de cisalhamento soldada sob a chapa de base, solução que evita
submeter os chumbadores a esforço combinado e que é usual em bases sujeitas a
força horizontal relevante. A chaveta, de 200 mm de comprimento e 80 mm de
embutimento, transfere esse cortante ao concreto por compressão localizada:
verificada como área parcialmente carregada, resulta em resistência de
266,67 kN e taxa de 0,596. Adotou-se fator de confinamento unitário, isto é, sem
o ganho que a norma permitiria, por depender a razão entre as áreas carregada e
de distribuição da geometria da sapata, fora do escopo.

A verificação do concreto é governada pelo cone de arrancamento, com 527,41 kN
de resistência de cálculo contra 1.781 kN da extração pela cabeça, resultando em
taxa de 0,526. Para que essa resistência se verifique, o bloco de fundação deve
garantir afastamento de borda não inferior a 600 mm em torno do grupo, o que
corresponde a dimensões aproximadas de 1,60 por 2,00 m em planta.

\subsection{Modos dispensados e limitação remanescente}\label{ssec:ancoragem-pendencias}

Dos demais modos de ruptura do concreto previstos na
EN 1992-4 \cite{EN1992-4:2018}, dois são dispensados por critério da própria
norma e o terceiro é verificado.

O arrancamento lateral é dispensado pelo item 7.2.1.9, por ser o afastamento de
borda de 600 mm superior a metade da profundidade efetiva, que vale 200 mm. A
ruptura de borda sob cortante é dispensada por não haver cortante nos
chumbadores, transmitido que é pela chaveta. O fendilhamento do concreto,
esse, é aplicável e foi verificado pela fórmula (7.26), com afastamento crítico
de 1.200 mm, resultando na mesma resistência de cálculo do cone, que permanece o
modo governante.

Resta uma limitação, e ela deve ser declarada. A ancoragem entre chumbador e
concreto não está validada por ensaio: os valores apresentados correspondem a
chumbador moldado com cabeça e decorrem de formulação normativa, não de
avaliação técnica do produto efetivamente especificado.

\section{CONTRAVENTAMENTO}\label{sec:detalhe-contraventamento}

O contraventamento da cobertura emprega cruzes de Santo André em barras redondas
de 28 mm em aço S235, dispostas apenas nas duas baías de extremidade, pelas
razões expostas na \autoref{ssec:predim-cv-cobertura}. O peso de aço do sistema
é de aproximadamente 486 kg.

O contraventamento vertical emprega oito painéis de cruz distribuídos pelas duas
fachadas longitudinais, com diagonais de 36 mm no mesmo aço. O contraventamento
transversal, nos frontões, emprega diagonais de 27 mm.

\subsection{Ligações de extremidade das diagonais}\label{ssec:lig-diagonais}

As diagonais terminam em olhal com pino, conforme o item 3.13 da
EN 1993-1-8 \cite{EN1993-1-8:2005}. Essas ligações não são dimensionadas para o
esforço solicitante, e sim por capacidade, isto é, para a resistência de cálculo
da própria barra. O critério assegura que, num evento que exceda a previsão, o
escoamento se dê na barra, que é dúctil e avisa, e não no olhal.

Na diagonal vertical, o esforço solicitante é de 166,88 kN e a resistência da
barra, de 211,69 kN, do que resulta sobre-resistência de 1,27. O
dimensionamento adota este último valor. A \autoref{tab:lig-diagonais} reúne as
duas ligações, conforme o item 3.13 e o item 4.5.3.3 da
EN 1993-1-8 \cite{EN1993-1-8:2005}, com chapa do olhal em aço S235 e pino em
S355. As taxas referem-se ao esforço de capacidade da barra, superior ao
solicitante.

\tabela
{Ligações de extremidade das diagonais de contraventamento}
{\footnotesize
\begin{tabular}{lcccc}
\toprule
    \textbf{Diagonal} & \textbf{Barra} & \textbf{Pino} & \textbf{Modo governante} & \textbf{$\eta$} \\
\midrule
    Vertical    & $\varnothing$ 36 mm & $\varnothing$ 36 mm & Geometria do olhal & 0,7713 \\
    Transversal & $\varnothing$ 27 mm & $\varnothing$ 27 mm & Geometria do olhal & 0,7633 \\
\bottomrule
\end{tabular}}
{Elaborado pelo autor (2026)}
{lig-diagonais}
{}
{}

A manilha e a rosca de extremidade são peças de catálogo e não são
dimensionadas neste trabalho.

%% TODO: dado ausente — especificar manilha e rosca com resistência de cálculo
%% à tração certificada pelo fabricante, não inferior a 211,7 kN na diagonal
%% vertical e a 119,1 kN na transversal.

A corda do diafragma merece registro construtivo. A terça de beiral acumula a
função de corda, recebendo esforço normal de 15,1 kN em ambos os sentidos, e sua
continuidade ao longo dos 46,0 m é condição para que o diafragma funcione. São
necessárias três emendas, e cada uma deve ser capaz de transmitir os 15,1 kN. A
omissão dessas emendas invalidaria o modelo de diafragma adotado, ainda que a
taxa de utilização da corda seja de apenas 0,0656.

\section{PROTEÇÃO CONSTRUTIVA}\label{sec:protecao-construtiva}

A classe de serviço 3 impõe exigências que não se resolvem por cálculo, e sim
por detalhe. O tratamento preservativo conforme a EN 351-1
\cite{EN351-1:2007} e o adesivo do tipo I conforme a EN 301 \cite{EN301:2017}
são condições necessárias, porém não suficientes.

O detalhamento deve assegurar que a água não permaneça em contato com a madeira.
As medidas pertinentes a esta edificação compreendem o afastamento das bases de
pilar em relação ao piso, por meio de pedestal de concreto que eleve a madeira
acima do nível de acúmulo de água; a proteção do topo das peças expostas; o
escoamento na calha central, que corre sobre a linha de pilares intermediários e
constitui o ponto mais crítico da cobertura; e a ventilação dos rasgos das
chapas embutidas, para que não retenham umidade.

\section{ESFORÇOS ENTREGUES À FUNDAÇÃO}\label{sec:cargas-fundacao}

Conforme delimitado na \autoref{sec:delimitacao}, a fundação não é dimensionada
neste trabalho. O programa encerra a cadeia de verificações entregando, por
base, os esforços de cálculo que constituem dado de entrada para o projeto
geotécnico. A \autoref{tab:cargas-fundacao} os reúne, nos valores de cálculo das
combinações governantes de cada efeito, conforme a EN 1990 \cite{EN1990:2002}. A
tração corresponde a arrancamento, sem desconto do peso próprio da sapata nem do
solo sobrejacente, e a força horizontal total transmitida é de 847,4 kN na
direção longitudinal e de 53,1 kN na transversal.

\tabela{Esforços de cálculo entregues ao projeto de fundação}{\footnotesize
\begin{tabular}{lcrrr}
\toprule
    \textbf{Base} & \textbf{Quantidade} & \textbf{Compressão} & \textbf{Tração} & \textbf{Horizontal} \\
    & & \textbf{(kN)} & \textbf{(kN)} & \textbf{(kN)} \\
\midrule
    Extremidade &  14 &  99,6 &  83,9 &   0,0 \\
    Central     &   9 & 199,2 & 151,8 &   0,0 \\
    Canto       &   4 & 293,2 & 277,5 & 159,0 \\
\bottomrule
\end{tabular}}{Elaborado pelo autor (2026)}{cargas-fundacao}{}{}

Não estão verificados, por dependerem de investigação geotécnica inexistente até
o momento, a tensão admissível do solo, as dimensões da sapata em planta, o
deslizamento, o tombamento, o recalque e a armadura. A única exigência
geométrica que este trabalho impõe ao bloco é a decorrente da ancoragem de
canto, exposta na \autoref{ssec:ancoragem-canto}.

\section{ESTIMATIVA PRELIMINAR DE CUSTO}\label{sec:estimativa-custo}

As quantidades apuradas ao longo dos capítulos anteriores permitem compor uma
estimativa de custo da estrutura. Trata-se de estimativa de anteprojeto,
destinada a situar a ordem de grandeza do investimento e a comparar
alternativas, e não de orçamento executivo: as parcelas que não decorrem
diretamente das quantidades calculadas são percentuais declarados, e os preços
unitários são valores de mercado que exigem cotação formal de fornecedor antes
de qualquer uso contratual.

\subsection{Quantidades}\label{ssec:custo-quantidades}

O volume de madeira lamelada colada distribui-se pelos três tipos de elemento
conforme a \autoref{tab:custo-quantidades}. O volume das vigas é o apurado na
\autoref{sec:geometria-adotada}, por integração da seção ao longo do vão. O dos
pilares decorre das seções e das quantidades da
\autoref{tab:geometria-edificacao} e da \autoref{tab:arrancamento}, com os
7,0 m de altura. O das terças decorre da seção de 0,20 por 0,58 m e do
espaçamento de 2,53 m aplicados à área de 1.863 m\textsuperscript{2} em planta,
relação que reproduz a parcela de 0,1889 kN/m\textsuperscript{2} registrada na
\autoref{tab:peso-proprio}.

\tabela{Quantidades de madeira lamelada colada e de aço de contraventamento}{\footnotesize
\begin{tabular}{llrr}
\toprule
    \textbf{Material} & \textbf{Elemento} & \textbf{Quantidade} & \textbf{Unidade} \\
\midrule
    MLC GL24h & Vigas principais (18)      & 220,8 & m\textsuperscript{3} \\
              & Terças                     &  85,4 & m\textsuperscript{3} \\
              & Pilares (27)               &  16,8 & m\textsuperscript{3} \\
\addlinespace
    \textbf{Subtotal MLC} & & \textbf{323,0} & \textbf{m\textsuperscript{3}} \\
\midrule
    Aço S235 & Contraventamento vertical, $\varnothing$ 36 mm    & 1.158,1 & kg \\
             & Contraventamento da cobertura, $\varnothing$ 28 mm &   485,8 & kg \\
             & Contraventamento transversal, $\varnothing$ 27 mm  &   162,9 & kg \\
\addlinespace
    \textbf{Subtotal aço} & & \textbf{1.806,8} & \textbf{kg} \\
\bottomrule
\end{tabular}}{Elaborado pelo autor (2026)}{custo-quantidades}{}{}

As massas de aço não são estimadas por percentual: decorrem do comprimento
total de barra de cada sistema, obtido da geometria dos painéis, multiplicado
pela área da seção circular e pela massa específica de 7.850 kg/m\textsuperscript{3}.
A massa do contraventamento da cobertura, de 485,8 kg, é a mesma registrada na
\autoref{sec:detalhe-contraventamento}. A do contraventamento vertical
corresponde ao diâmetro de 36 mm da execução consolidada, e não ao de 33 mm
registrado na \autoref{tab:varredura-paineis}, cuja divergência já consta
daquela seção.

\subsection{Preços unitários adotados}\label{ssec:custo-unitarios}

O mercado brasileiro pratica dois patamares distintos para a madeira lamelada
colada. O primeiro, entre R\$ 2.000 e R\$ 2.800 por metro cúbico, corresponde à
peça trabalhada em oficina, sem montagem, e representa o custo de produção
interno de um fabricante. O segundo, entre R\$ 4.000 e R\$ 6.400 por metro
cúbico, corresponde ao preço de venda comercial do produto estrutural acabado,
já incluída a margem do fabricante.

Adotou-se o segundo patamar, com valor central de R\$ 5.200 por metro cúbico.
A razão é que a estimativa se coloca na posição do projetista ou da construtora
que adquire a estrutura, e não na de quem a fabrica: o custo relevante é o de
aquisição, não o de produção. O intervalo completo é examinado adiante, na
\autoref{ssec:custo-sensibilidade}.

Duas ressalvas incidem sobre esse valor. A primeira é que os preços de mercado
reunidos referem-se a setembro de 2026 e variam com a região, o fornecedor e o
volume contratado, de modo que a estimativa envelhece. A segunda é de natureza
técnica: a classe GL24h especificada no \autoref{cap:definicao-problema} decorre
da EN 14080 \cite{EN14080:2013}, ao passo que os preços praticados no Brasil
correspondem a madeira lamelada colada de pinus ou de eucalipto, cuja
certificação segue outra referência. A equivalência entre ambas é hipótese de
orçamento assumida aqui, e não resultado do dimensionamento.

%% TODO: referência — inserir em referencias.bib a fonte dos preços de mercado
%% da MLC (gerador de preços de construção civil ou tabela de fabricante
%% nacional), com data de coleta, e substituir a menção genérica por \cite{}.

%% TODO: dado ausente — o preço unitário do aço S235 em barra redonda,
%% fabricada e montada, de R$ 18,00/kg, é valor declarado de ordem de grandeza e
%% não decorre de cotação. Confirmar com fornecedor antes da entrega.

%% TODO: verificar — confirmar junto a fabricante nacional de MLC se a classe
%% GL24h da EN 14080 é fornecida no Brasil, e a que preço, ou identificar a
%% classe equivalente na ABNT NBR 7190-1 e o seu preço de mercado.

\subsection{Composição do custo}\label{ssec:custo-composicao}

A \autoref{tab:custo-composicao} compõe o custo em cinco parcelas. As duas
primeiras decorrem das quantidades apuradas e dos preços unitários. A terceira,
relativa às chapas, pinos e chumbadores das ligações, adota quinze por cento do
custo da madeira, percentual de ordem usual em estruturas de madeira lamelada
colada com ligações por chapa metálica embutida, e não foi apurada peça a peça.
A quarta cobre a mão de obra de montagem, que não está incluída no preço de
venda do produto acabado, e adota doze por cento do subtotal de materiais. A
quinta corresponde aos \criarsigla*{BDI}{benefícios e despesas indiretas},
adotados em vinte e cinco por cento, dentro da faixa usual de orçamento de
edificações.

\tabela{Composição do custo estimado da estrutura}{\footnotesize
\begin{tabular}{llr}
\toprule
    \textbf{Parcela} & \textbf{Critério} & \textbf{Valor (R\$)} \\
\midrule
    Material MLC        & 323,0 m\textsuperscript{3} $\times$ R\$ 5.200/m\textsuperscript{3} & 1.679.600 \\
    Aço de contraventamento & 1.806,8 kg $\times$ R\$ 18,00/kg                                &    32.522 \\
    Ligações e ferragens    & 15\% do material MLC                                            &   251.940 \\
    Montagem                & 12\% do subtotal de materiais                                   &   235.687 \\
\addlinespace
    \textbf{Subtotal direto} & & \textbf{2.199.750} \\
\midrule
    BDI                     & 25\% do subtotal direto                                         &   549.937 \\
\addlinespace
    \textbf{Total estimado} & & \textbf{2.749.687} \\
\bottomrule
\end{tabular}}{Elaborado pelo autor (2026)}{custo-composicao}{}{}

O total corresponde a R\$ 1.476 por metro quadrado de área coberta, tomada a
projeção em planta de 1.863 m\textsuperscript{2}. A madeira responde por
sessenta e um por cento do total, e o aço de contraventamento, apesar de
governar o dimensionamento com a taxa de 0,7902 da
\autoref{tab:verif-contraventamento}, responde por pouco mais de um por cento.

%% TODO: verificar — os percentuais de 15\% para ligações, 12\% para montagem e
%% 25\% para BDI são valores declarados de prática usual e não decorrem de
%% composição de custo unitário. Confirmar com referência de orçamento (CUB,
%% CBIC ou tabela de composições) e registrar a fonte em referencias.bib.

\subsection{Sensibilidade ao preço unitário}\label{ssec:custo-sensibilidade}

A incerteza dominante não está nas quantidades, que decorrem do dimensionamento,
e sim no preço unitário da madeira. A \autoref{tab:custo-sensibilidade} percorre
o intervalo de mercado, mantidas as quantidades e os percentuais.

\tabela{Custo total estimado em função do preço unitário da MLC}{\footnotesize
\begin{tabular}{lrrr}
\toprule
    \textbf{Cenário} & \textbf{R\$/m\textsuperscript{3}} & \textbf{Total (R\$)} & \textbf{R\$/m\textsuperscript{2}} \\
\midrule
    Limite inferior de mercado & 4.000 & 2.125.651 & 1.141 \\
    \textbf{Adotado}           & \textbf{5.200} & \textbf{2.749.687} & \textbf{1.476} \\
    Limite superior de mercado & 6.400 & 3.373.723 & 1.811 \\
\bottomrule
\end{tabular}}{Elaborado pelo autor (2026)}{custo-sensibilidade}{}{}

A amplitude é de cinquenta e nove por cento entre os extremos, o que dimensiona
a precisão desta estimativa: ela distingue ordens de grandeza e ordena
alternativas, mas não substitui cotação. Uma consequência de projeto decorre
dessa faixa. A varredura da \autoref{sec:varredura-global} identificou
geometrias aprovadas com cerca de metade do volume de madeira da adotada, o que
corresponde, aos preços aqui empregados, a uma diferença da ordem de meio
milhão de reais no material. As razões que sustentaram a geometria mais robusta,
expostas na \autoref{ssec:o-que-governa}, são de margem nos sistemas de
estabilização e de construtibilidade, e é contra essa ordem de grandeza que
devem ser pesadas.
